\begin{table}[ht]
\centering
\caption{Acurácia e ganho relativo por perturbação ($T=0{,}6$). (*) em \textbf{Ganho}: variação significativa entre Original e a perturbação no mesmo modelo (McNemar, $p<0{,}05$).}
\label{tab:acuracia_completa}
\sisetup{output-decimal-marker={,}, table-format=-2.1, table-space-text-post={*}}
\resizebox{\textwidth}{!}{%
\begin{tabular}{l
  S[table-format=2.1, table-space-text-post={*}]
  S[table-format=-2.1, table-space-text-post={*}]
}
\toprule
\textbf{Perturbação} & \multicolumn{1}{c}{\textbf{Acurácia (\%)}} & \multicolumn{1}{c}{\textbf{Ganho (\%)}} \\
\cmidrule(lr){2-2} \cmidrule(lr){3-3}
& {LRM} & {LRM} \\
\midrule
Original & 57.0 & \text{--} \\
Linguagem Natural & 63.9 & +12.1* \\
Embaralhado & 59.3 & +4.0 \\
Junto & 58.2 & +2.0 \\
Irrelevante & 58.0 & +1.7 \\
Informação Ausente & 55.8 & -2.0* \\
Complexo & 56.8 & -0.3 \\
Contradição & 46.8 & -17.8* \\
Negação & 17.7 & -69.0* \\
\bottomrule
\end{tabular}%
}
\end{table}